\documentclass[11pt,a4paper]{article}
\usepackage[margin=22mm]{geometry}
\usepackage{amsmath,amssymb,amsthm,booktabs,microtype,needspace}
\usepackage[T1]{fontenc}
\usepackage{lmodern}
\usepackage[colorlinks=true,linkcolor=blue,urlcolor=blue]{hyperref}
\hypersetup{pdftitle={Subset sums: an integer prefix inverse and the audited bound 93},pdfauthor={Sum-of-Subset research project}}
\newcommand{\FS}{\operatorname{FS}}
\newcommand{\Z}{\mathbb Z}
\newcommand{\T}{\mathcal T}
\newtheorem{theorem}{Theorem}[section]
\newtheorem{lemma}[theorem]{Lemma}
\newcommand{\RootGroups}{11,052}
\newcommand{\CoarseGroups}{10,958}
\newcommand{\NestedGroups}{94}
\newcommand{\IntervalChecks}{11,090}
\newcommand{\ExtremumPoints}{23,917}
\newcommand{\TreeRoots}{284}
\newcommand{\BinaryNodes}{943}
\newcommand{\TreeLeaves}{1,227}
\newcommand{\MaxDepth}{12}
\newcommand{\FinalSystems}{76,891}
\newcommand{\Vectors}{1,539}
\newcommand{\MinGap}{1/196}
\title{Subset Sums in a Fourfold Sumset\\\large An integer prefix inverse and the audited bound $N\le93$}
\author{Sum-of-Subset research project}
\date{16 October 2026}
\begin{document}
\maketitle
\begin{abstract}
We continue the independently audited chain from $N\le94$ to $N\le93$.
A new exact integer inverse combines the threshold-17 non-progression
lemma with a dichotomy on large excess and non-progression prefixes.
We use small-hole witnesses from the first admissible parameter, rather
than waiting until the hole count reaches 15. At cardinality 94 this
reduces the positive-low-hole covering cap from 266 to 242, below the
three-empty threshold 258. The low-hole root universe shrinks by 20,264
integer systems. The high-hole proof retains the newly possible two-empty
family and groups 3,178 states into 69 common models while retaining their hole variable.
Fresh exact certificates close the complete endpoint; archived endpoint
certificates are not used. The stronger archived bound 84 has not yet
received a complete independent replay, so this is progress in the
audited chain, not a numerical improvement over that archive.
\end{abstract}

\section{Result and inherited foundation}
For a finite real set $A$, let
$\FS(A)=\{\sum_{x\in X}x:X\subseteq A\}$ include zero, and let
$4A=A+A+A+A$ permit repeated summands. Let $N$ be the maximum feasible
size under $\FS(A)\subseteq4A$.
\begin{theorem}\label{main}
Every feasible set has at most 93 elements. Hence $71\le N\le93$.
\end{theorem}
The lower bound is still
$[-31,31]\cap\Z\cup\{\pm93,\pm155,\pm217,\pm713\}$.
The upper bound follows from the audited chain
$115\to103\to100\to99\to98\to97\to96\to95\to94$ and the exclusion of
size 94 below. We inherit the corrected peeling, integralization,
small-increment inverse, simultaneous inverse, effective-hole, covering,
mass, span, and sharp four-extreme lemmas from
\cite{foundation,overlap,previous}. Their hypotheses are checked anew.
No feasibility under arbitrary deletion is assumed.

For a positive $r$-set $X$, put $T_r=r(r+1)/2$ and
\[
 \mu(X)=|\FS(X)|-1-T_r,\quad
 \delta_0(X)=|\{x\in X:2x\notin R(X\cup\{0\})\}|,\quad
 \rho(X)=\mu(X)+2-2\delta_0(X),
\]
where $R$ consists of distinct-element two-sums. Increasing prefixes
$X_i$ have nonnegative integer increments
$y_i=|\FS(X_i)|-|\FS(X_{i-1})|-i$, with
\begin{equation}\label{rec}
 \mu(X_r)=\sum_i y_i,\qquad
 \rho(X_i)\ge\rho(X_{i-1})+y_i-2.
\end{equation}
The previous report proves $\rho(X)\ge r-5$ for $r\ge17$ and
non-progression $\FS(X)$, with sharp constant 5. Its 2,723 finite
base cases are recomputed by the new verifier, without modifying that
lemma or its independent checker.

\section{An exact non-progression prefix inverse}
The inherited inverse, for $2\le\beta\le10$, is
\[
 L_\beta(r,u)=\min\{(\beta-1)(r-14),
 (\beta-1)q+\min(s,\beta-1)\},
 \quad \max(0,u-4\beta+2)=(\beta+1)q+s,
\]
where $0\le s\le\beta$. For a non-progression set of size $r\ge17$ with
$\mu\ge u$, define the valid old floor
\begin{equation}\label{old}
 V(r,u)=\max\{r-5,u+2-2r,L_2(r,u),\ldots,L_{10}(r,u)\}.
\end{equation}
The number 14 inside $L_\beta$ is retained: substituting 5 there would
not follow from the sharper non-progression theorem.

For integers $v\ge0$ and $k\ge1$, write $v=(k+2)q+s$, $0\le s<k+2$,
and define
\begin{equation}\label{inverse}
 J_k(v)=kq+\min(s,k).
\end{equation}
This is the least nonnegative integer $t$ with
$t+2\lfloor t/k\rfloor\ge v$. Indeed, on a block $t=kq+h$ with
$0\le h<k$, the left side is $(k+2)q+h$; the two skipped values
are first reached at the next multiple of $k$.

\begin{theorem}[Prefix inverse]\label{prefix}
Let $r\ge17$, let $\FS(X)$ be non-progression, and suppose $\mu(X)\ge u$.
For $17\le L\le r$ and $3\le\beta\le\min(10,L-4)$, put
\begin{align*}
 R_0&=r-5,&M&=T_L-r,\\
 A&=R_0+(\beta-2)(r-L),&C&=V(r,M+1),\\
 B&=R_0+J_{\beta-2}\bigl(\max(0,u-R_0-4\beta+2)\bigr).
\end{align*}
Then $\rho(X)\ge\min(A,B,C)$. Consequently
\begin{equation}\label{new}
 F(r,u)=\max\left(V(r,u),\ \max_{L,\beta}\min(A,B,C)\right)
\end{equation}
is a valid lower bound for $\rho(X)$.
\end{theorem}
\begin{proof}
If $\mu\ge M+1$, use \eqref{old} to obtain $\rho\ge C$.
Otherwise every prefix of size at least $L$ has non-progression support.
To see this, if a prefix at index $i\ge L$ were a progression, let
$k>i$ be the first subsequent non-progression index. Its predecessor
support and new translate must be disjoint: overlap puts the translating
element on the same step and their union is another progression.
It follows that
\[
 y_k=|\FS(X_{k-1})|-k\ge1+T_{k-1}-k\ge1+T_L-r=M+1,
\]
contradicting $\mu=\sum y_i\le M$.

If every $y_i>\beta$ for $L<i\le r$, the floor at $X_L$ and
\eqref{rec} give $\rho\ge L-5+(\beta-1)(r-L)=A$.
Otherwise take the last $i>L$ with $y_i\le\beta$. The inverse theorem
applies because $i\ge L+1\ge\beta+5$. Up to scaling, $X_i$ is an
interval with at most $\beta$ deletions, retaining its maximum.
It has at most $2\beta$ diagonal defects. Here the elementary
endpoint-band argument is sufficient: for $q\ge2$ deletions and $i>q$,
every defect lies in $[1,q]$ or $[i,i+q]$, and not all $q$ low elements
can be present and defective, since defectivity of 1 excludes 2.
Thus there are at most $(q-1)+(q+1)=2q$ defects. For $q=1$ there
are at most two, and for $q=0$ only the maximum.

Write $t=r-i$. Each later step adds at most one defect and raises
$\rho$ by at least $\beta-1$, so
\[
 \delta_0(X)\le2\beta+t,\qquad
 \rho(X)\ge i-5+(\beta-1)t=R_0+(\beta-2)t.
\]
With $s=\rho-R_0\ge0$, the identity for $\mu$ yields
\[
 u-R_0-4\beta+2\le s+2\left\lfloor\frac{s}{\beta-2}\right\rfloor.
\]
The exact inverse \eqref{inverse} gives $\rho\ge B$.
These cases exhaust all configurations and prove the theorem.
\end{proof}
The archive \cite{archive} contains a related prefix argument starting
at 21 with a linear ceiling. The present proof starts at 17, permits
$\beta\le L-4$ by searching strictly after $L$, and preserves the integer
staircase rather than replacing the floor by a real quotient.
For example, $(r,u)=(38,165)$ gives $V=105$ but $F=111$, witnessed by
$(L,\beta,M,A,B,C)=(20,8,172,141,111,112)$.
The checker verifies the exact inverse at 5,409 parameter pairs
$1\le k\le9$, $0\le v\le600$, including minimality. This is a check
of the finite formula used here; the proof of the theorem is the
unbounded argument above.

\section{The size-94 lattice gate}
A size-94 counterexample contains zero, since adjoining zero otherwise
would give a size-95 counterexample. After integralization, remove four
original extremes on each side to get $D$, and peel twice to
$D\supset D_1\supset E=P\cup\{0\}\cup(-Q)$. Use
\begin{gather*}
 d=86,\quad e=78,\quad m=a+b=77,\quad I=(a-38)(a-39),\\
 \Lambda=1+T_a+T_b=1522+I,\quad\Omega=467-I.
\end{gather*}
The inherited identity is
\[
 \rho(P)+\rho(Q)+2\vartheta+c+\T=\Omega,
 \qquad c=2|2D|+18d+85-|\FS(D)|,
\]
with nonnegative terms, and $\T=\kappa+\ell_1+\ell_2+g$.
The deficient-sign test is $1+T_{78}>24\cdot94+77$.
The paired bootstrap gives $q^2\le2477$ with odd $q$, so $q\le49$
and $a,b\ge14$. Then $\Omega\ge0$ leaves $17\le a\le60$.
The simultaneous inverse condition holds for all 44 orientations:
\[
 T_{\min(a,b)}-2-\Omega\ge38>0.
\]
It supplies a common core lattice, with holes
$z=h_P+h_Q-\kappa$, $0\le z\le\min(h_P,h_Q)\le\kappa$ and
$h_X<|\FS(X)|-1$.

Write $\rho=\rho(P)+\rho(Q)$, $\mu=\mu(P)+\mu(Q)$ and
$C_0=2m+4e+14-\Lambda=-1042-I$. The inherited inner comparisons are
\[
 |\FS(D_1)|-2|\FS(E)|\le
 \min\{C_0+\Omega+\rho-\mu+2z,\ C_0+3\Omega-\rho-\mu\}.
\]
For $z\le14$, $\rho\le\mu$ makes the first bound at most $-547-2I$.
For $z\ge15$, both supports are non-progressions. The size-17 theorem
already applies to the cores, without assuming the outer sets are on
the lattice, giving $\rho+\mu\ge(m-10)+(m-6)=138$.
Whenever this is insufficient, use the effective-hole witness at
$\beta=7$ separately for each side. The necessary rows are
\begin{center}\small
\begin{tabular}{lrrrrrrr}\toprule
$r$&32&33&34&35&36&37&38\\
$C_7(r)$&83&89&96&102&109&110&116\\\midrule
$r$&39&40&41&42&43&44&45\\
$C_7(r)$&122&128&129&135&142&147&149\\\bottomrule
\end{tabular}
\end{center}
They give $\mu(X)\ge r+C_7(r)+1$. With $\rho\ge m-10$, the final
high-hole comparison is at most $-3$. The outer margin is
\[
 \Lambda-(4d-10+\Omega)=721+2I>0.
\]
Thus an off-lattice element would produce disjoint support translates,
contradicting the relevant strict factor-two comparison. Both $D_1$ and
$D$ are on the core lattice. Only now use hole propagation and the sharp
four-extreme theorem to obtain
\begin{equation}\label{budget}
 w\le4z,\qquad\T\ge\kappa+162-3z,\qquad
 \T\ge\max(54,z,162-2z).
\end{equation}
The sharp non-progression theorem is valid for real positive sets and
is independent of \eqref{budget}; this order of proof is not circular.

\section{Hole-sensitive pruning and coupled groups}
We use the effective-hole theorem from the foundation in its exact form.
Given positive integers $r,\beta$ and $C\ge0$, let $K$ be the least
integer $\ge3$ with $(K+1)(K-2)/2>C+1$, and put
\begin{gather*}
 D_0=C+K,\quad E_0=D_0+1,\\
 M_0=\max\left(2\beta,1+\left\lfloor
 \frac{1+\sqrt{9+8(E_0+6\beta)}}2\right\rfloor\right),\quad
 J=\max(K+1,2\beta+5,M_0).
\end{gather*}
A witness is accepted only when all three integer inequalities hold:
\[
 T_{M_0}-6\beta>M_0+1+E_0,\qquad
 \beta(r-J+1)>D_0,\qquad T_{J-1}-1>D_0.
\]
If $2\beta<h_X<|\FS(X)|-1$, it then gives $\mu(X)\ge r+C+1$.
The search uses $17\le r\le60$, $1\le\beta\le20$, $0\le C\le600$;
only valid witnesses are used, and no optimality of these limits is
claimed. Every chosen witness is checked again with exact integers.

For $z>0$, let $U(r,z)$ be the maximum of $r-3$ and all bounds
$r+C+1$ from accepted witnesses with $2\beta<z$.
For $z=1,2$ only $r-3$ is needed. Set $R(r,z)=F(r,U(r,z))$.
The common covering cap is
\begin{equation}\label{cap}
 c\le467-I-R(a,z)-R(b,z)-\max(54,z,162-2z).
\end{equation}
For $z=0$ use zero side penalties, obtaining $c\le305-I$.
All positive-hole states are enumerated within $1\le z\le\Omega$;
negative caps are discarded. The finite scalar audit gives
\[
 1\le z\le14\ \Longrightarrow\ c\le242-I<3d.
\]
For $z=1,2$ use $R(a,z)+R(b,z)\ge67$. For $z=3,4$, only the
$\beta=1$ witnesses are needed: \eqref{new} gives a combined floor
at least 80 over all 44 orientations. For $5\le z\le14$, the
$\beta=2$ witnesses alone give at least 103. Thus
\[
 c+I\le\max\{305-67+4,\ 305-80+8,\ 305-103+28\}=242.
\]
All 88 pairs of small-hole witnesses and inverse parameters are saved
in the scalar audit. This step uses hole information before any transport.
The former coarse estimate would only give $c\le266-I$.

Let $j$ be the number of empty covering units and $\delta=c-86j$.
The complete low-hole root ranges are
\begin{center}\small
\begin{tabular}{lrrrrrrr}\toprule
Regime&$Z_0$&$Z_1$&$Z_2$&$Z_3$&$P_0$&$P_1$&$P_2$\\\midrule
$a$&22--55&24--53&27--50&32--45&23--54&26--51&31--46\\
Integer roots&136&2,100&13,872&41,888&128&1,820&9,248\\\bottomrule
\end{tabular}
\end{center}
Here $Z$ means zero holes and $P$ means positive low holes.
The old cap 266 would retain 31,460 positive roots, including 17,952
three-empty roots. The new positive count is 11,196: 20,264 fewer roots
before geometric solving. With the unchanged 57,996 zero-hole roots,
the complete low-hole universe has 69,192 integer models.

For $z\ge15$, 3,896 oriented scalar states survive and $c\le184<3d$.
The first pair's actual contributor fixes the direction in the inherited
active-center estimate; reflection cannot be chosen to maximize it.
There remain 1,350 states for $j=0$, 1,746 for $j=1$, and 82 for $j=2$.
For each fixed $(a,j,U(a,z),U(b,z))$, retain the actual hole count as
one common variable $z$, bounded by the smallest and largest values in
the group. The side excess floors are constant throughout this group.
Write $E_*=467-I-F(a,U(a,z))-F(b,U(b,z))$. Instead of independently
maximizing the hole count and the cap, keep all three coupled rows
\[
 \delta+z\le E_*-86j,\qquad
 \delta\le E_*-54-86j,\qquad
 \delta-2z\le E_*-162-86j.
\]
They are exactly the three alternatives inside the maximum in
\eqref{cap}. They remain correlated with the same $z$ in the mass,
span and cardinality rows. The common model may include unused integer
hole values and real intermediate values; this is a valid relaxation,
not an assertion that those values arise from a set. The additional
constant cap $\delta\le C_*-86j$ is retained to select the row policy.
For every actual state the checker verifies membership in the interval,
the common side floors, the equality of the scalar and three-row caps,
and domination by the group model, with the group's fixed row policy.
There are 974,444 row checks over all 3,178 retained states.
For example, the group $(a,b,u_p,u_q)=(36,41,148,193)$ has
$c\le177$ and $z\le145$, but also $c+z\le231$.
The pair of separate maxima $(177,145)$ violates this last inequality
by 91. It is a relaxation artifact, not a feasible set.
There are 17 zero-empty, 46 one-empty and 6 two-empty groups, hence
\[
 1350\cdot4+1746\cdot70+82\cdot578=175016
 \quad\longrightarrow\quad17\cdot4+46\cdot70+6\cdot578=6756
\]
high-hole root models. For positive low holes, use the established
transport to $z=14$ after imposing the sharper, $z$-independent cap.

\section{Complete necessary models and conditional rows}
We retain the twenty-center, eighteen-unit covering grammar and all
actual-contributor choices from \cite{overlap,previous}. The numbers of
modes with $j=0,1,2,3$ empty units are $4,70,578,2992$.
Missing units may not be removed from the root list. In the potentially
split pair, either actual member is retained as a separate choice;
its positivity inequality is necessary, not sufficient for contribution.

Write $x_0=R_p<x_1<x_2<x_3<x_4=M_p$, and similarly $y_i$ on the other
side, and put $K=|2D|$. The unchanged mass rows on side $p$ are
\begin{align*}
 2M_p+6R_p-2S_p+2z&\le-2T_a-12-2u_p,\\
 M_p+4R_p-S_p&\le-T_{a-1}-6,\\
 x_1+x_2+x_3+M_p-S_p+z&\le-T_a-u_p,\\
 R_p+x_1+x_2+x_3+M_p-S_p&\le-T_{a-1},
\end{align*}
with $u_p=0$ for $Z$, $u_p=a-3$ for $P$, and $u_p=U(a,z)$ (constant in the group) in a
high group. The gap rows include $R_p\ge a$,
$M_p-R_p\ge4$ and $x_i-x_{i-1}\ge1$. Reflection supplies side $q$.
The new global numerical rows are
\begin{align*}
 \delta&\le305-86j-I &&(Z),\\
 \delta&\le242-86j-I &&(P),\\
 R_p+R_q+\delta-4z&\le386-I-86j,\\
 S_p+S_q+\delta-2K-4z&\le1632-86j.
\end{align*}
For high groups use the three coupled cap rows and retain the same variable
$z$ in the span and cardinality inequalities; only the low-hole positive models
are transported to $z=14$.

The row policy is determined from a proved group-wide residual cap
$C_g$, never supplied by the certificate. If $r_{\min}$ is the smallest
side size in the group, the conditions are
\begin{center}\small
\begin{tabular}{ll}\toprule
Rows&Required condition\\\midrule
Ordinary quantiles&$C_g<r_{\min}+5$\\
Core quantiles&$C_g<82$\\
Basic leading quantiles&$C_g<171$\\
Nested leading quantiles&$C_g<163$\\
Strict outer order&$j\le1$, or $j=2$ with $C_g<171$\\\bottomrule
\end{tabular}
\end{center}
Choose the strongest justified individual quantile type. At a positive
center $v$ the global inequalities are
$v+l+M_p\ge86$ and $v-u-M_q\le S_p-85$; core versions replace
$M_p,M_q,86,85$ by $R_p,R_q,82,81$. Ordinary versions are
$v+l\ge b+5$ and $v-u\le S_p-a-4$.
Individual loss caps are 85 globally and 81 for core rows.
In particular, $Z_3$ now has cap 47 and uses core rows, not the
unconditional ordinary rows valid at the preceding endpoint.
High two-empty groups have cap at most 12 and satisfy the strict-outer
condition. The high one-empty cap can be 98, so core rows are not
universally justified there.

Joint block variables represent the actual whole-block covering deficit.
For $k$ ordinary centers in a block the valid bound is
\[
 O\ge78k-(R_p+R_q+v_{\max}-v_{\min}+1)+\max_v(l_v+u_v).
\]
Let $V$ be the retained own centers in a block. For its leading positive
slot and a nonempty subset $S\subseteq V$,
\[
 O\ge78(|S|+2)-2-p_{\max}-R_p-2R_q
       +\max\{L_p,l_v+u_v:v\in S\}
       +\sum_{v\in V\setminus S}(l_v+u_v).
\]
The analogous sum-of-individual-loss bounds also hold. Each is a
lower bound on the same $O$; the separate gains are not added together.
For an own-center ordinary block, its leading-slot loss is added to the
ordinary hull and individual-loss bounds, as in the inherited grouping proof.
Only disjoint blocks are summed, with ungrouped losses, to at most
$\delta$. These are the proved grouping inequalities of \cite{overlap},
now with $e=78$. Leading-center quantiles use $2d-1=171$ and
$2e+7+2i=163+2i$, respectively, under their stated cap tests.

The coarse model substitutes $K\le M_p+M_q+R_p+R_q+21$ into the
cardinality row and eliminates $K$. The nested model retains all five
interval-envelope levels and the 25 independent-cutoff tail rows
\[
 K\le M_p+M_q+x_i+y_h+1+T_{4-i}+T_{4-h},\qquad0\le i,h\le4.
\]
All models use nonnegative real coordinates as a relaxation of actual
configurations. Four proved binary alternatives are available: core
quantiles or loss at least 82; ordinary quantiles or an outside endpoint;
a loss count at least $i$ or its inner-element cutoff; and both
contributors or a sole contributor in the split pair. Both children
must be certified. Endpoint alternatives after transport are chosen
using the enlarged endpoint, exactly as proved in \cite{previous}.
No bound on the total number of large-loss centers is imposed.

\section{Exact certificate coverage and reproduction}
The verifier rebuilds the scalar domain, all mode choices, row policies,
and the complete root list before reading certificate assignments.
There are $69192+6756=75948$ integer root systems. The final proof has
\begin{center}\small
\begin{tabular}{lr}\toprule
Quantity&Count\\\midrule
Structural groups&\RootGroups\\
Coarse / nested groups&\CoarseGroups\ / \NestedGroups\\
Symbolic interval checks&\IntervalChecks\\
Integer extremum evaluations&\ExtremumPoints\\
Integer roots requiring trees&\TreeRoots\\
Binary internal nodes / leaves&\BinaryNodes\ / \TreeLeaves\\
Maximum tree depth&\MaxDepth\\
Final integer leaf systems&\FinalSystems\\
Distinct integer vectors&\Vectors\\\bottomrule
\end{tabular}
\end{center}
For $Ax\le b$, $x\ge0$, each nonnegative integer certificate vector $w$
satisfies $w^TA\ge0$ and $w^Tb<0$, giving a strict contradiction.
The smallest normalized gap $-w^Tb/\sum w_i$ is $\MinGap$.
Symbolic interval checking uses directly constructed quadratic right-hand
sides, exact rational endpoints, and the integers around the vertex;
it never infers a polynomial from samples. Every integer parameter is
assigned once. Binary leaves are checked individually, and omitted
children, unknown splits, repeated splits and unused vectors are rejected.

The 23 negative controls include omitted high two-empty modes, the
invalid universal core shortcut, arbitrary reflection, an unproved
size-16 prefix inverse, a wrong integer inverse, independent hole/cap maxima, and an attempted
exclusion of size 93. The conclusion is $N\le93$, not $N\le92$.
\clearpage
\subsection*{Reproduction and scope}
Only Python 3.10+ and its standard library are required:
\begin{verbatim}
python -I -S -B verification/nonap_inverse.py
python -I -S -B verification/upper93.py --output .work/upper93.json
python -I -S -B verification/upper93.py --exhaustive
\end{verbatim}
The data file is \path{data/upper93_certificates.json}.
The exhaustive mode rebuilds every integer model, checks every row
against its symbolic expression, and checks each multiplier pointwise.
Here the expanded integer parameter is $a$. In a high-hole group, $z$
remains a model coordinate: each exact row combination certifies every
real value in its interval simultaneously, including every actual integer
hole count, rather than testing only the endpoints of that interval.
The optional \path{experiments/generate_upper93_certificates.py} uses
SciPy only for discovery; every saved certificate must pass exact checking.
Results, the nested-group index, generation logs, report review and input
hashes are in \path{results/2026-10-16/upper93/}.

Standard and optimized Python outputs agree byte for byte. Exhaustive
and interval results agree apart from the mode label. Regeneration is
byte-identical, and upper-94 regression is unchanged.
\path{verify_upper93_windows.bat} runs all nine finite continuations.
The infinite mathematical reductions are proved here and in the cited
foundation; program PASS is not a proof-assistant formalization.
The next endpoint is cardinality 93 and needs its own lattice and
complete coverage checks before any further bound can be asserted.

\begin{thebibliography}{9}\raggedright
\bibitem{previous} Sum-of-Subset research project,
\emph{A sharper non-progression lemma and the audited bound $N\le94$},
\path{reports/upper94_from95_20261016.tex}, 16 October 2026.
\bibitem{overlap} Sum-of-Subset research project,
\emph{Joint covering losses and the audited bound $N\le95$},
\path{reports/upper95_from96_20261016.tex}, 16 October 2026.
\bibitem{foundation} Sum-of-Subset research project,
\emph{The audited continuation from 115 to 103},
\path{reports/upper103_from115_20261014.tex}, 14 October 2026.
\bibitem{archive} Sum-of-Subset research project,
\emph{Complete upper-86 research report},
\path{reports/upper86_complete_20261005.tex}, 5 October 2026.
The subsection ``Dense spans and the stronger non-AP inverse'' motivates
Section 2; no archived endpoint exclusion is assumed.
\end{thebibliography}
\end{document}
